% ======================================================================
% Tensorial strain reconstruction from eddy-current testing in LPBF.
% elsarticle, lean structure (<=4 sections), standalone TikZ/pgfplots figures.
% Build:  make    (figures first, then this file)   --- or latexmk -pdf main
% NOTE: author affiliations are the authors' institutions; confirm before submission.
% ======================================================================
\documentclass[preprint]{elsarticle}
% Submission:  \documentclass[review]{elsarticle}
% Production:  \documentclass[final,5p,times,twocolumn]{elsarticle}

\usepackage{amsmath,amssymb,bm}
\usepackage{graphicx}
\usepackage{xcolor}
\usepackage{booktabs}
\usepackage{siunitx}
\usepackage{hyperref}
\usepackage[capitalise]{cleveref}   % after hyperref

\graphicspath{{figures/}}

\newcommand{\TODO}[1]{\textcolor{red}{[TODO: #1]}}
\newcommand{\dsig}{\Delta\sigma}
\newcommand{\eps}{\bm{\varepsilon}}
\newcommand{\sig}{\bm{\sigma}}
\newcommand{\Kt}{\mathsf{K}}

\journal{\TODO{NDT\&E International / Additive Manufacturing}}

\begin{document}

\begin{frontmatter}

\title{Tensorial strain reconstruction from eddy-current testing in
laser powder bed fusion}

\author[pv]{M.\ Carraturo\corref{cor}}
\ead{massimo.carraturo@unipv.it}
\cortext[cor]{Corresponding author}
\author[tum]{B.\,I.\ Ates}

\affiliation[pv]{organization={Department of Civil Engineering and Architecture,
University of Pavia}, city={Pavia}, country={Italy}}
\affiliation[tum]{organization={Technical University of Munich},
city={Munich}, country={Germany}}

\begin{abstract}
Eddy-current testing (ECT) is an attractive in-situ probe for laser powder bed
fusion (LPBF): non-contact, fast, and sensitive to the electrical conductivity,
which depends on mechanical strain. A calibrated scalar law relates the relative
conductivity change to strain, but the residual strain in an as-built part is a
second-order tensor whereas a coil reports a single scalar, so reconstructing
six components from one measurement is under-determined. We present a framework
that combines (i) a tensorial elastoresistivity model generalising the scalar
coefficient to a fourth-order coupling tensor; (ii) a finite-element (FEM)
reciprocity analysis quantifying which conductivity components a probe senses;
(iii) optimal design of directional-probe sets (a conductivity rosette); and
(iv) a regularised, prior-informed maximum-a-posteriori inversion that fuses the
measurements with a process simulation. An axisymmetric eddy-current FEM,
verified to better than $10^{-5}$ by an ohmic-energy balance, agrees with the
analytical model to \SI{0.9}{\percent} in resistance and \SI{2.8}{\percent} in
reactance. On synthetic LPBF builds a normal coil resolves one strain degree of
freedom, an in-plane rosette three, and a tilted six-probe set all six; the
prior supplies the unobserved components. The workflow is realised in an open,
tested pipeline that outputs a full strain-tensor field.
\end{abstract}

\begin{keyword}
eddy-current testing \sep laser powder bed fusion \sep residual strain \sep
elastoresistivity \sep inverse problems \sep in-situ monitoring
\end{keyword}

\end{frontmatter}

% \linenumbers

% ======================================================================
\section{Introduction}\label{sec:intro}
% ======================================================================
Residual stress and strain are central concerns in metal additive manufacturing:
they drive distortion, cracking, and premature failure, and they are a defining
challenge of laser powder bed fusion (LPBF)~\cite{bartlett2019}. In-situ
monitoring able to map the evolving strain field during a build would enable
closed-loop process control, yet residual strain is still measured only
post-mortem, by destructive or diffraction methods~\cite{withers2001}.
Eddy-current testing (ECT) is a promising in-situ modality~\cite{garcia2011}: a
coil driven at frequency $\omega$ induces eddy currents in the conductive part,
and the change in coil impedance encodes the local electrical conductivity
$\sigma$. Because $\sigma$ depends on elastic strain through the elastoresistive
(piezoresistive) effect, ECT senses strain indirectly.

A calibrated scalar law,
\begin{equation}
  \frac{\dsig}{\sigma_0} = \kappa\,\varepsilon,
  \qquad \kappa = -0.326 \quad(\text{316L}),
  \label{eq:scalar}
\end{equation}
relates the relative conductivity change to a uniaxial strain $\varepsilon$, with
coefficient of determination $R^2 = 0.992$ in our uniaxial-tension calibration of
316L stainless steel. Equation~\eqref{eq:scalar} is, however, a one-dimensional
projection of a richer reality: in three dimensions the strain $\eps$ is a
symmetric second-order tensor with six independent components, while a coil
reports a single scalar impedance. Recovering six numbers from one is impossible
without additional information.

That conductivity depends on elastic strain underpins an established line of
eddy-current residual-stress assessment: high-frequency apparent-conductivity
spectroscopy and iterative inversion recover near-surface residual-stress depth
profiles in shot-peened superalloys~\cite{yu2004,abunabah2006,abunabah2007}. That
work also exposed a key confounder --- cold work and dislocation-density changes
alter conductivity alongside elastic strain --- which is especially pronounced in
LPBF microstructures and which we revisit in \cref{sec:results}. In metal
additive manufacturing, recoater-mounted ECT is already used for layer-wise
density and defect monitoring~\cite{spurek2022,monu2025}, with part-temperature
compensation shown to be essential~\cite{spurek2024}; strain, however, is not yet
mapped in~situ.

This paper develops the methodology to close that gap. Our contributions are: a
tensorial elastoresistivity model that generalises $\kappa$ to a fourth-order
coupling tensor; an FEM-based observability analysis, built on a verified and
cross-validated forward model, that quantifies which conductivity components a
probe senses; an optimal directional-probe (``rosette'') design; and a
prior-informed tensor inversion realised in an open pipeline. The proposed
deployment is sketched in \cref{fig:overview}.

\begin{figure}[tb]
  \centering
  \includegraphics[width=\linewidth]{fig_overview}
  \caption{Proposed deployment, doubling as a graphical abstract. A
  recoater-mounted, multi-frequency directional-probe rosette scans each layer;
  the directional measurements $y=\hat{\bm n}^{\!\top}\sig\,\hat{\bm n}$ are fused
  with a thermomechanical process-simulation prior by a regularised (MAP)
  inversion to reconstruct the strain-tensor field $\varepsilon_{ij}(\bm x)$.}
  \label{fig:overview}
\end{figure}

% ======================================================================
\section{A tensorial eddy-current framework for strain reconstruction}
\label{sec:method}
% ======================================================================

\subsection{Eddy-current forward model}\label{sec:forward}
\paragraph{Analytical model.} For a cylindrical multi-turn coil of inner/outer
radius $r_i,r_o$, length $\ell$, $n$ turns and lift-off $s$ above a conductive
half-space, the impedance follows the Dodd--Deeds
formulation~\cite{dodd1968,bowler2019,theodoulidis2006}
\begin{equation}
  Z = \frac{\mathrm{j}\omega\mu_0\pi n^2}{\ell^2 (r_o-r_i)^2}
  \int_0^\infty \frac{J^2(\alpha r_i,\alpha r_o)}{\alpha^5}
  \Big[\,2\ell + \tfrac{1}{\alpha}\big(2e^{-\alpha\ell}-2 + P(\alpha)\,\Gamma(\alpha)\big)\Big]\,
  \mathrm{d}\alpha,
  \label{eq:dd}
\end{equation}
with $P(\alpha)=e^{-2\alpha(\ell+s)}+e^{-2\alpha s}-2e^{-\alpha(\ell+2s)}$,
the kernel $J(\alpha r_i,\alpha r_o)=\int_{\alpha r_i}^{\alpha r_o}xJ_1(x)\,
\mathrm{d}x$, and $\mu=\mu_0$ for the non-magnetic 316L considered here. The
half-space reflection coefficient is
\begin{equation}
  \Gamma(\alpha)=\frac{\alpha-\gamma}{\alpha+\gamma},
  \qquad \gamma=\sqrt{\alpha^2+\mathrm{j}\omega\mu\sigma},
  \label{eq:refl}
\end{equation}
which tends to $-1$ as $\sigma\to\infty$, so that eddy currents reduce the net
inductance and add resistance. The normalised impedance is
$Z_{\mathrm{norm}}=(Z_0-Z)/\operatorname{Im}(Z_0)$, with $Z_0$ the free-space
(self) impedance.

\paragraph{Axisymmetric finite-element model.} To treat geometries where the
half-space assumption fails (finite parts, layers, defects), we solve the
time-harmonic eddy-current problem for the azimuthal magnetic vector potential
$A_\phi$ in the meridional $(\rho,z)$ plane (\cref{fig:domain}). Starting from
$\nabla\times(\mu^{-1}\nabla\times\bm A)+\mathrm{j}\omega\sigma\bm A=\bm J_s$ with
$\bm A=A_\phi\,\hat{\bm\phi}$, the symmetric weak form reads
\begin{equation}
  \frac{1}{\mu}\!\int\!\Big[\nabla A_\phi\!\cdot\!\nabla v
  + \frac{A_\phi v}{\rho^2}\Big]\rho\,\mathrm{d}\rho\,\mathrm{d}z
  + \mathrm{j}\omega\!\int\!\sigma A_\phi v\,\rho\,\mathrm{d}\rho\,\mathrm{d}z
  = \int J_s\, v\,\rho\,\mathrm{d}\rho\,\mathrm{d}z
  \label{eq:weak}
\end{equation}
for all test functions $v$, with the Dirichlet condition $A_\phi=0$ on the
symmetry axis ($\rho=0$, where the azimuthal component must vanish) and on the
truncated outer boundary. Imposing these conditions is essential: without them
the $1/\rho^2$ term renders the system singular. The domain comprises four
regions (\cref{fig:domain}): a conductor bulk, a near-surface conductor
\emph{skin layer} meshed to one-third of the skin depth
$\delta=\sqrt{2/(\omega\mu\sigma)}$ to resolve the exponential eddy-current decay,
an air gap, and the air above, with the winding modelled as a uniform azimuthal
current density $J_s=nI/S_{\mathrm{coil}}$ over its rectangular cross-section
$S_{\mathrm{coil}}=(r_o-r_i)\ell$. We discretise with $H^1$ Lagrange elements of
order three and solve the complex system directly.

The coil impedance is obtained from the magnetic flux linkage of the winding,
\begin{equation}
  \Lambda = \frac{n}{S_{\mathrm{coil}}}\!\iint_{\mathrm{coil}}\! 2\pi\rho\,A_\phi\,
  \mathrm{d}\rho\,\mathrm{d}z, \qquad Z=\mathrm{j}\omega\Lambda/I,
  \label{eq:flux}
\end{equation}
and $Z_0$ is computed by the \emph{same} model with $\sigma=0$, so that any
discretisation or truncation error cancels consistently in the normalisation.
\Cref{tab:params} lists the reference parameters.

\begin{figure}[tb]
  \centering
  \includegraphics{fig_domain}
  \caption{Axisymmetric FEM computational domain in the meridional $(\rho,z)$
  plane. The conductor carries a near-surface skin layer meshed to $\sim\delta/3$;
  the coil winding cross-section carries a uniform azimuthal current density
  $J_s$; $A_\phi=0$ holds on the symmetry axis and the truncated outer boundary.
  Lift-off $s$, coil dimensions $r_i,r_o,\ell$ and skin depth $\delta$ are marked
  (not to scale).}
  \label{fig:domain}
\end{figure}

\begin{table}[tb]
  \centering
  \caption{Reference coil, material and FEM parameters.}
  \label{tab:params}
  \begin{tabular}{llll}
    \toprule
    Quantity & Symbol & Value & \\
    \midrule
    Coil inner/outer radius & $r_i,r_o$ & $4.04,\,11.84$ & \si{\milli\metre}\\
    Coil length / turns     & $\ell,n$  & $8.02$ \si{\milli\metre}, $1858$ & \\
    Lift-off                & $s$       & $1.0$ & \si{\milli\metre}\\
    Frequency               & $f$       & $240$ & \si{\kilo\hertz}\\
    Conductivity (reference)& $\sigma$  & $16.2$ & \si{\mega\siemens\per\metre}\\
    Domain (radial/axial)   & $R$; $z$  & $50$; $[-20,30]$ & \si{\milli\metre}\\
    Skin-layer mesh size    & $h$       & $\delta/3$ & \\
    Element order / DOF     & $p$       & $3$; ${\sim}4\times10^4$ & \\
    \bottomrule
  \end{tabular}
\end{table}

\subsection{Tensorial elastoresistivity}\label{sec:tensor}
Strain breaks the isotropy of the material, so under load the conductivity is a
second-order tensor $\sigma_{ij}$ that couples to strain through a fourth-order
\emph{elastoresistivity} tensor $\Kt$,
\begin{equation}
  \frac{\dsig_{ij}}{\sigma_0} = K_{ijkl}\,\varepsilon_{kl}.
  \label{eq:tensor}
\end{equation}
For an isotropic base material $\Kt$ takes the form
$K_{ijkl}=a\,\delta_{ij}\delta_{kl}+b(\delta_{ik}\delta_{jl}+\delta_{il}\delta_{jk})$
--- two independent constants. Writing them as a longitudinal $\kappa_\parallel$
and a transverse $\kappa_\perp$ coefficient,
\begin{equation}
  \frac{\dsig_{ii}}{\sigma_0}=\kappa_\parallel\varepsilon_{ii}
  + \kappa_\perp\!\!\sum_{k\neq i}\varepsilon_{kk}, \qquad
  \frac{\dsig_{ij}}{\sigma_0}=(\kappa_\parallel-\kappa_\perp)\varepsilon_{ij}
  \;\; (i\neq j).
  \label{eq:coupling}
\end{equation}
The scalar law~\eqref{eq:scalar} is the uniaxial projection of \eqref{eq:tensor}:
the relative conductivity change sensed along a unit direction $\hat{\bm n}$ in a
uniaxial-stress test is $\hat{\bm n}^{\!\top}(\dsig/\sigma_0)\hat{\bm n}$, which
varies strongly with $\hat{\bm n}$ (\cref{fig:rosette}). A single uniaxial test
fixes only one projection; we anchor $\kappa_\parallel$ so that the along-load
projection reproduces the measured scalar
($\kappa_\parallel-2\nu\kappa_\perp=-0.326$ at $\nu=0.30$), giving
$\kappa_\parallel=-0.374$ for an illustrative $\kappa_\perp=-0.08$. Crucially, the
observability results depend on these constants only through the ratio
$\kappa_\perp/\kappa_\parallel$ (\cref{sec:results}); separating the two requires
multi-orientation loading. A purely volumetric response
($\kappa_\parallel=\kappa_\perp$) makes $\Kt$ rank-one, so conductivity then
constrains only the volumetric strain $\operatorname{tr}\eps$.

\subsection{Observability and probe design}\label{sec:obs}
By driving-point reciprocity~\cite{auld1999}, a small conductivity perturbation
changes the coil impedance by
\begin{equation}
  \delta Z = -\frac{1}{I^2}\int \delta\sigma\, \bm E\!\cdot\!\bm E \,\mathrm{d}V .
  \label{eq:recip}
\end{equation}
For the axisymmetric coil the field is purely azimuthal,
$\bm E=-\mathrm{j}\omega A_\phi\,\hat{\bm\phi}$, so $\bm E\!\cdot\!\bm E=E_\phi^2$
and only the local in-plane (azimuthal) component of a tensor perturbation is
sensed: a normal coil is structurally \emph{blind} to $\sigma_{zz}$, and its
sensitivity is confined to within a skin depth of the surface.

A directionally selective probe senses, to leading order, the effective
conductivity along its sensing direction,
\begin{equation}
  y(\hat{\bm n}) = \hat{\bm n}^{\!\top}\sig\,\hat{\bm n},
  \label{eq:proj}
\end{equation}
the conductivity analogue of a mechanical strain-gauge rosette.
Equation~\eqref{eq:proj} is a leading-order model: it assumes current flows
cleanly along $\hat{\bm n}$. Accessed from the build surface this is well
justified in-plane but optimistic out-of-plane, since eddy currents flow
essentially parallel to the surface --- the axisymmetric result above makes this
exact. A coil merely tilted above the surface therefore gains only limited
$\sigma_{zz}$ sensitivity; genuine out-of-plane access requires edge/side
probing, elongated or rotating-field probes, or multi-frequency depth profiling.
We accordingly treat the full six-direction result as an \emph{upper bound} on
directional diversity, and take the in-plane rosette combined with the simulation
prior (\cref{sec:method}, \cref{eq:map}) as the practical configuration.

Stacking $M$ probes gives a linear map
$\bm y=\mathsf{M}\,\bm\sigma_{\mathrm v}$ from the conductivity Voigt vector to the
measurements; composing with $\Kt$ yields the strain map
$\bm y=\mathsf{A}\,\eps_{\mathrm v}$, $\mathsf{A}=\sigma_0\,\mathsf{M}\,\Kt$. The
singular value decomposition of $\mathsf A$ gives the resolvable rank, the
conditioning, and the null space (blind directions). We design probe sets by
minimising the condition number of $\mathsf A$ subject to a rank constraint; the
resolvable ceiling is $\operatorname{rank}\Kt$, so directional diversity helps
only insofar as the elastoresistive response is anisotropic. Excitation frequency
adds a complementary, depth-selective degree of freedom through $\delta$.

\subsection{Prior-informed tensor inversion and pipeline}\label{sec:inv}
Because realistic probe sets are rank-deficient, we estimate strain by a
maximum-a-posteriori (MAP) inversion~\cite{aster2018} that fuses the measurements
with a process-simulation prior $\eps_{\mathrm{prior}}$ (data assimilation). With
Gaussian measurement noise ($\sigma_d$) and prior ($\sigma_p$),
\begin{equation}
  \eps^\star = \eps_{\mathrm{prior}}
  + \big(\mathsf A^{\!\top}\mathsf A + \lambda \mathsf I\big)^{-1}
  \mathsf A^{\!\top}\big(\bm y - \mathsf A\,\eps_{\mathrm{prior}}\big),
  \qquad \lambda=\sigma_d^2/\sigma_p^2 .
  \label{eq:map}
\end{equation}
The resolution matrix
$\mathsf R=(\mathsf A^{\!\top}\mathsf A+\lambda\mathsf I)^{-1}\mathsf A^{\!\top}
\mathsf A$ quantifies, per component, the fraction supplied by the data: its
diagonal is $1$ where a component is measurement-determined and $0$ where it is
taken from the prior. Where only a normal-coil scalar is available, the estimate
degrades gracefully to the volumetric strain via the invariant coupling
$\kappa_v=(\kappa_\parallel+2\kappa_\perp)/3$. The forward model, observability,
probe design and inversion are combined in a single open pipeline that ingests
layer-wise ECT data and outputs both the reconstructed geometry and, from
directional measurements, a full strain-tensor field; the implementation is
covered by an automated test suite. The scalar $\lambda$ assumes an
independent-and-identically-distributed prior and equal probe noise; a full prior
covariance from an ensemble of process simulations is a natural refinement.

% ======================================================================
\section{Results and discussion}\label{sec:results}
% ======================================================================
\paragraph{Forward-model verification.} The FEM is mesh-converged (order-three
and order-four solutions agree to five significant figures) and internally
consistent: the resistance from the flux-linkage impedance~\eqref{eq:flux} equals
a direct volume integral of the ohmic loss $\tfrac12\!\int\sigma|E|^2\mathrm{d}V$
to better than $10^{-5}$. Solving for the reference coil over a conductor of
$\sigma=\SI{16.2}{\mega\siemens\per\metre}$ at \SI{240}{\kilo\hertz} gives a
purely inductive air impedance $Z_0=\mathrm{j}\,5.18\times10^4~\Omega$ and a
loaded impedance $Z=622+\mathrm{j}\,3.73\times10^4~\Omega$ --- a positive
resistance (eddy-current loss) and a reduced reactance, as required. The
analytical and FEM normalised impedances agree to \SI{0.9}{\percent} in
resistance and \SI{2.8}{\percent} in reactance (\cref{fig:impedance}), the small
reactance residual being the finite FEM truncation versus the analytical
half-space. The analytical integral~\eqref{eq:dd} is converged for
$k_{\max}\!\sim\!6000$, within a stable plateau before the oscillatory Bessel
kernel degrades the quadrature, and the reciprocity kernel~\eqref{eq:recip} is
confirmed against finite-difference re-solves, the ratio approaching unity as the
perturbation shrinks (\cref{fig:femconv}). Maintaining two dissimilar forward
models is what made these checks possible: an inverse method that re-uses a
single model for synthesis and inversion cannot detect its own forward-model
error.

\begin{figure}[tb]
  \centering
  \includegraphics{fig_impedance}
  \caption{Normalised-impedance locus traced by the analytical model as the
  conductivity sweeps $\sigma\in[10^4,10^8]~\si{\siemens\per\metre}$, with the FEM
  result at the reference conductivity. The two agree to \SI{0.9}{\percent} in
  resistance and \SI{2.8}{\percent} in reactance.}
  \label{fig:impedance}
\end{figure}

\begin{figure}[tb]
  \centering
  \includegraphics[width=\linewidth]{fig_fem_convergence}
  \caption{Forward-model verification. Left: the analytical normalised impedance
  is stable over the integration cut-off $k_{\max}$. Right: the reciprocity
  sensitivity kernel matches finite-difference re-solves, the ratio
  $\to 1$ as the conductivity perturbation $\delta\sigma/\sigma_0\to 0$.}
  \label{fig:femconv}
\end{figure}

\paragraph{What a probe senses.} The FEM sensitivity density $|E_\phi|^2$ in the
specimen decays from the surface over a skin depth (\cref{fig:kernel}), and
excitation frequency tunes that penetration: \SI{120}{\kilo\hertz} reaches
markedly deeper than \SI{480}{\kilo\hertz} (\cref{fig:depth}). Because the field
is in-plane, a single coil senses only the in-plane mean conductivity; the
effective $\kappa$ it reports depends on the sensing direction, swinging from
$-0.35$ along the load to $+0.06$ transverse (\cref{fig:rosette}). Stacking
probes and analysing $\mathsf A$ then shows the resolving power: a normal coil
resolves one strain degree of freedom, an in-plane rosette three, and a tilted
six-direction set all six (\cref{fig:spectra}, \cref{tab:obs}). The tilted-set
result is an upper bound presuming idealised out-of-plane access.

\begin{figure}[tb]
  \centering
  \includegraphics{fig_kernel}
  \caption{FEM sensitivity density $|E_\phi|^2$ in the specimen $(\rho,z)$ plane,
  normalised to its maximum. The measurement is confined to within a skin depth
  of the surface and a coil radius of the axis.}
  \label{fig:kernel}
\end{figure}

\begin{figure}[tb]
  \centering
  \includegraphics{fig_depth}
  \caption{Mean in-specimen sensitivity versus depth at two frequencies; higher
  frequency concentrates the sensitivity nearer the surface, providing a
  depth-selective degree of freedom.}
  \label{fig:depth}
\end{figure}

\begin{figure}[tb]
  \centering
  \includegraphics{fig_rosette}
  \caption{Conductivity ``rosette'': the effective scalar $\kappa$ sensed in a
  uniaxial-stress test as a function of the in-plane sensing angle. A single
  scalar reading is therefore ambiguous; directional diversity is required.}
  \label{fig:rosette}
\end{figure}

\begin{figure}[tb]
  \centering
  \includegraphics{fig_spectra}
  \caption{Singular-value spectra of the strain map $\mathsf A$ for three probe
  configurations. The number of non-negligible singular values equals the
  resolvable strain degrees of freedom: one, three and six respectively.}
  \label{fig:spectra}
\end{figure}

\begin{table}[tb]
  \centering
  \caption{Resolvable strain degrees of freedom and conditioning by probe set.
  The tilted six-probe row presumes idealised out-of-plane access and is an upper
  bound.}
  \label{tab:obs}
  \begin{tabular}{lcc}
    \toprule
    Configuration & resolvable DOF (of 6) & condition number\\
    \midrule
    Single normal coil            & 1 & 1.0\\
    In-plane rosette ($0/45/90$)  & 3 & 2.57\\
    In-plane rosette ($0/60/120$) & 3 & 2.18\\
    Optimised in-plane rosette    & 3 & 1.93\\
    Rosette $+$ tilted probes (6) & 6 & 2.50\\
    \bottomrule
  \end{tabular}
\end{table}

\paragraph{Material limit on observability.} The resolvable ceiling is
$\operatorname{rank}\Kt$, set by the material. Sweeping the ratio
$\kappa_\perp/\kappa_\parallel$ (\cref{fig:kappa}) shows that the resolvable DOF
stay full for any anisotropic response, but the conditioning of $\mathsf A$
degrades smoothly toward the volumetric limit $\kappa_\perp\!\to\!\kappa_\parallel$,
where the rank collapses to one. At the calibration-anchored ratio
($\approx0.21$) both the rosette and the full set are well conditioned, so the
design is insensitive to the under-determined split between $\kappa_\parallel$ and
$\kappa_\perp$.

\begin{figure}[tb]
  \centering
  \includegraphics[width=\linewidth]{fig_kappa}
  \caption{Strain observability versus the elastoresistivity ratio
  $\kappa_\perp/\kappa_\parallel$. Left: condition number of $\mathsf A$, which
  diverges as the response becomes volumetric. Right: resolvable DOF, full for any
  anisotropic response. The dotted line marks the calibration-anchored ratio.}
  \label{fig:kappa}
\end{figure}

\paragraph{Tensor reconstruction.} On a synthetic build-height profile with a
deliberately biased prior, the inversion~\eqref{eq:map} corrects the in-plane
components from data under both probe sets, while the out-of-plane components are
recovered only by the full set and otherwise follow the prior
(\cref{fig:inversion}) --- exactly as the resolution diagonal predicts
(\cref{fig:resolution}). The full set reduces the strain RMSE from
$9.4\times10^{-4}$ (prior only) to $2.6\times10^{-4}$, and the in-plane rosette to
$3.9\times10^{-4}$. Embedded in the pipeline, the method reconstructs a full
strain-tensor field; \cref{fig:fieldeqv,fig:fieldvol} show the derived
von-Mises-equivalent and volumetric strain fields for a $60\times40$ build, the
volumetric strain evolving from compressive near the base to tensile higher in
the build. Each voxel is currently inverted independently; the finite footprint
of the sensitivity kernel (\cref{fig:kernel}) couples neighbours and should, in a
quantitative treatment, be folded into $\mathsf A$.

\begin{figure}[tb]
  \centering
  \includegraphics[width=\linewidth]{fig_inversion}
  \caption{Tensor-strain reconstruction along the build height with a biased
  prior. The in-plane $\varepsilon_{xx}$ is corrected from data by both
  configurations; the out-of-plane $\varepsilon_{zz}$ is recovered only by the
  full set and otherwise follows the prior.}
  \label{fig:inversion}
\end{figure}

\begin{figure}[tb]
  \centering
  \includegraphics{fig_resolution}
  \caption{Data-resolved fraction (resolution-matrix diagonal) per strain
  component. The in-plane rosette resolves $xx,yy,xy$ from data but leans on the
  prior for $zz$ and the out-of-plane shears.}
  \label{fig:resolution}
\end{figure}

\begin{figure}[tb]
  \centering
  \includegraphics{fig_field_eqv}
  \caption{Reconstructed von-Mises-equivalent strain field over the build,
  recovered by the full probe set.}
  \label{fig:fieldeqv}
\end{figure}

\begin{figure}[tb]
  \centering
  \includegraphics{fig_field_vol}
  \caption{Reconstructed volumetric strain $\operatorname{tr}\eps$, from
  compressive (blue) near the base to tensile (orange) higher in the build.}
  \label{fig:fieldvol}
\end{figure}

\paragraph{Feasibility on LPBF 316L.} Three physical realities set the instrument
requirements. (i) \emph{Signal:} residual elastic strains of
$1$--$2\times10^{-3}$ give a relative conductivity signal of only
$\dsig/\sigma_0\approx3$--$7\times10^{-4}$. (ii) \emph{Temperature, the dominant
confounder:} 316L has $\rho\approx\SI{0.74}{\micro\ohm\metre}$
($\sigma\approx\SI{1.35}{\mega\siemens\per\metre}$) with a resistivity temperature
coefficient of order $10^{-3}~\si{\per\kelvin}$~\cite{nist316l}, so a drift of
${\sim}\SI{1}{\kelvin}$ changes $\sigma$ by as much as the entire strain signal
--- consistent with in-situ ECT studies that found temperature dominant and
required compensation~\cite{spurek2024}. Differential or reference-channel probes,
multi-frequency lift-off/temperature separation, and tight thermal referencing are
therefore essential; cold work and dislocation density are a further,
microstructure-borne confounder. (iii) \emph{Depth:} at the realistic 316L
conductivity the skin depth is ${\approx}\SI{0.9}{\milli\metre}$ at
\SI{240}{\kilo\hertz} (and ${\approx}\SI{0.4}{\milli\metre}$ at
\SI{1}{\mega\hertz}), far larger than a $30$--$\SI{50}{\micro\metre}$ layer, so
each measurement averages $20$--$30$ layers; layer-resolved tomography would
require multi-frequency depth inversion, and the reference-coil verification above
(a higher, generic conductivity) should be read as a verification case rather than
a 316L prediction. These constraints do not undermine the framework --- they
specify the instrument: a thermally-referenced, multi-frequency
directional-probe rosette.

% ======================================================================
\section{Conclusions}\label{sec:conclusions}
% ======================================================================
We have shown how to progress from a scalar strain--conductivity relation to a
tensorial reconstruction with ECT. The scalar impedance constrains a
\emph{projection} of the conductivity tensor; resolving the strain tensor then
requires measurement diversity (a directional, optionally multi-frequency
rosette) and, for the components that remain unobservable from the surface, a
physics-based prior fused by regularised inversion. A verified axisymmetric FEM
agrees with the analytical model to \SI{0.9}{\percent} in resistance and
\SI{2.8}{\percent} in reactance; a normal coil resolves one strain degree of
freedom, an in-plane rosette three, and a tilted six-probe set up to six; and a
MAP inversion reduces the strain RMSE on a synthetic build from
$9.4\times10^{-4}$ (prior only) to $2.6\times10^{-4}$ with the full set. The main
limitations are out-of-plane observability (surface ECT accesses in-plane current
paths far more readily than $\sigma_{zz}$, which leans on the prior or
multi-frequency profiling) and calibration (multi-orientation tests are needed to
separate $\kappa_\parallel$ and $\kappa_\perp$ and to characterise the as-built
$\sigma_0$ anisotropy). Future work targets a three-dimensional anisotropic
eddy-current model to quantify the true out-of-plane sensitivity of tilted and
elongated probes, experimental calibration against multi-orientation and
benchmark~\cite{burke1988} data, and deployment of the pipeline for in-situ
monitoring.

% ----------------------------------------------------------------------
\appendix
\section{Reference implementation}\label{app:impl}
The forward models, observability and probe-design tools, and the prior-informed
inversion are implemented in an open Python library with an automated test suite;
every figure in this paper is regenerated from version-controlled data by the
scripts in the repository (see \emph{Code and data availability}).

% ----------------------------------------------------------------------
\section*{Code and data availability}
The open-source pipeline and the scripts and data that regenerate every figure
are available at \url{https://github.com/MassimoCarraturo/eddy_current_project};
an archived release with a DOI will be deposited on Zenodo.

\section*{CRediT author statement}
\textbf{M.\ Carraturo:} conceptualisation, methodology, software, writing ---
original draft. \textbf{B.\,I.\ Ates:} methodology, calibration, writing ---
review and editing.

\section*{Declaration of generative AI}
Generative-AI tools assisted software development and manuscript editing; all
scientific content was verified by the authors, who take full responsibility for
it.

\bibliographystyle{elsarticle-num}
\bibliography{refs}

\end{document}
